\chapter{Amendment Log}
\label{back:amendments}

\section*{Amendments Ratified Since the Founding}

The following amendments have been ratified by the Assembly of the silo since this Pact was first adopted. Each is entered below in the order of its ratification, with the text of the amendment and a note of its effect. The Pact does not itself reckon a calendar from the founding, and this Log follows the Pact's own practice: an amendment's place here is fixed by its order, not by a year that binds it to a count no citizen is asked to keep.

\subsection*{Amendment 1: The Housing Proviso}

\textit{Ratified in the silo's early years, not long after the offices of this Pact had settled into their ordinary practice.}

``Section 1 of Article 14 is amended to add the following clause: `A citizen whose dwelling is revoked under this Section may petition Judicial to reduce or stay the revocation, provided the citizen has made a good-faith effort to remedy the breach that occasioned it and sufficient time has passed for remedy to have been possible.' The remainder of Section 1 is unchanged.''

\textit{Effect}: Judicial is given explicit discretion to allow remediation before imposing housing revocation. This amendment reflects the principle that citizens are given opportunity to correct small breaches before suffering larger penalties.

\subsection*{Amendment 2: The Testimony Entered into the Lottery}

\textit{Ratified in the generation after Amendment 1, upon a petition of citizens rather than the Mayor's proposal, the first amendment so brought.}

``Article 3, Section 6 is amended to add, before its clauses, the following: `It is recorded, though not by the name of the household, that a citizen who waited three years for a lottery slot, and buried a stillborn child in the fourth, came before the Assembly and said that this balance, hard as it had been to bear, was the same balance that had kept her own mother fed as a child, and asked that no citizen be spared the waiting who had not first been spared the alternative. This Pact does not require any citizen to feel as she felt. It asks only that every citizen know such a citizen has lived, and judged as she judged.' The remainder of Section 6 is unchanged.''

\textit{Effect}: No requirement of the Section is altered; a citizen's testimony is set within it, that the lottery's necessity be heard in a citizen's voice as well as the Pact's. The Assembly's record notes that the petitioners first asked the Judge whether a citizen's words might be entered into the Pact at all, and that the Judge ruled they might, provided the household went unnamed, as no other citizen is named anywhere in the Pact.

\subsection*{Amendment 3: The Rumor Proviso}

\textit{Ratified in the same generation as Amendment 2, upon the Mayor's proposal, after a season the Mayor's own report to the Assembly called a season of loose talk.}

``Article 16, Section 2 is amended to add: `Forasmuch as a rumor is a leak in the hull no less than a crack in the wall, and a hull is kept whole by the many hands that stop it and not by the one that made it, every citizen who hears a rumor of the kind this Section forbids shall report it to the Sheriff within one hour of the hearing, and the Sheriff shall enter the report, the hour, and the name of the citizen reporting, and no report so entered shall be held against the citizen who made it, whatever the rumor proved.' The remainder of Section 2 is unchanged.''

\textit{Effect}: A duty to report rumor is added to the duty Section 1 already lays upon a citizen who hears the forbidden word, and an hour is set for it. The Judge afterward ruled, under Article 6, Section 2, that the hour is reckoned from the citizen's reaching a landing where a Deputy is posted, and not from the hearing itself, no citizen upon the stair being able to keep it otherwise; the ruling stands entered beside the amendment.

\subsection*{Amendment 4: The Retirement Entitlement}

\textit{Ratified some years after Amendment 2.}

``Article 19, Section 5 is amended to require that a retired citizen receive, in addition to the ordinary ration, a medicine stipend adequate to the cost of common treatments, and that the Infirmary provide the same care to a retired citizen as to a citizen still in labor. The remainder of Article 19 is unchanged.''

\textit{Effect}: Retired citizens are guaranteed medication and continued medical care rather than left to purchase such care from their ration chits.

\subsection*{Amendment 5: The Appeal of Judicial Sentence}

\textit{Ratified some years further on.}

``Article 6, Section 4 is amended to add: `A citizen sentenced to the mines under Article 17 may petition Judicial to review the sentence once the citizen has served half the term originally set, and the Judge may shorten the remainder of the term if the citizen has demonstrated good conduct and honest labor throughout the sentence's duration, or if material new evidence has emerged bearing on the citizen's guilt. No such review shall reduce a sentence's total length by more than half of what was originally imposed.' The remainder is unchanged.''

\textit{Effect}: Citizens in the mines are given a pathway to early release at the midpoint of good behavior, though no sentence may be halved more than once, and the core of every sentence must still be served.

\subsection*{Amendment 6: The Right to Counsel}

\textit{Ratified in the years following Amendment 5.}

``Article 6, Section 4 is amended to add: `A citizen accused of an offense not among the gravest named under Article 17 may request that another citizen be permitted to speak on the accused's behalf before Judicial, that the Judicial hearing not be conducted in isolation. Judicial shall grant such request, provided the speaking-citizen is not themselves under prosecution and has not been barred from Assembly by prior judgment.' The remainder is unchanged.''

\textit{Effect}: Citizens are permitted advocates in non-gravest trials, though the form and scope of advocacy remain within Judicial's discretion.

\subsection*{Amendment 7: The Bounding of the Forgiveness Holiday}

\textit{Ratified some years after Amendment 6, upon the Mayor's own proposal, in a season the Mayor's report called a season of many surrenders.}

``Article 4, Section 7 is amended to add: `No Forgiveness Holiday shall be declared within five years of the ending of the last, nor shall any run longer than one turn of the season; and the Mayor shall report to the Assembly, at the sitting next following, the number of things surrendered and the number classified among the gravest, without the names of the citizens who surrendered them.' The remainder of Section 7 is unchanged.''

\textit{Effect}: The Holiday, which the Pact left wholly to the Mayor's judgment as to interval and cause, is bounded in both, and a count of what it yields is made known to the Assembly.

\subsection*{Amendment 8: The Syndromic Retirement}

\textit{Ratified well after Amendment 6, in a season this Log does not further mark.}

``Article 13, Section 8 is amended to add: `A citizen diagnosed with the Syndrome may petition for retirement under the same process as any other citizen seeking to leave labor, without waiting for the age at which retirement is ordinarily granted. The Head Physician and the Mayor shall approve such petition within one cycle of its filing.' The remainder is unchanged.''

\textit{Effect}: Citizens with the Syndrome are permitted to seek retirement at any age if their condition makes continued labor untenable.

\subsection*{Amendment 9: The Restoration of the Holiday}

\textit{Ratified after Amendment 8, upon a petition of citizens, the second amendment so brought. The text as ratified does not survive in the archive entire: the leaf of the roll bearing it is wanting, and what stands here is the Judge's reconstruction from the archive's marginal entry and the recollection of the tellers, entered as authoritative by a ruling under Article 6, Section 2, at the sitting next following.}

``Article 4, Section 7, as amended by Amendment 7, is restored to its form before that amendment, Amendment 7 being repealed entire [\ldots] save that the count of things surrendered [\ldots] be kept [\ldots] as the Mayor judges [\ldots].''

\textit{Effect}: The bound upon the Holiday is lifted, and the Mayor's judgment as to its interval and cause stands as the Pact first gave it. Whether the count of things surrendered is still required, the reconstruction does not settle; the Mayor's reports since have given it in some years and not in others.

\subsection*{Amendment 10: The Floor Reassignment Procedure}

\textit{Ratified some time after Amendment 8.}

``Article 14, Section 2 is amended such that: `A citizen may petition the Head of Supply for reassignment to a different floor's housing, and the Head of Supply shall consider such petition with regard to the citizen's trade, the silo's current needs, and the space available on the requested floor. A denial of such petition may be appealed to Judicial, though Judicial shall confirm the decision where the Head of Supply's reason is plainly stated and grounded in the silo's need.' The remainder is unchanged.''

\textit{Effect}: Citizens gain a formal right to petition for floor reassignment, with Judicial review available for denials.

\subsection*{Amendment 11: The Assembly Quorum}

\textit{Ratified later still.}

``Article 20, Section 1 is amended: `An extraordinary Assembly session to address a matter of urgent silo need may proceed with not fewer than one third of the silo's population of majority present in the hall or counted from the floors under Section 6, provided the session has been called in accordance with this Article and all citizens have been given not fewer than seven days' notice; a session with fewer present may hear, but may not vote.' The remainder is unchanged.''

\textit{Effect}: The Pact had set no least number for an extraordinary session; this amendment sets one, that a vote taken with few present not be afterward disputed. The ordinary Assembly, to which any citizen of majority may come or send word under this Article, is unaffected.

\subsection*{Amendment 12: The Reporting of Anomalous Signal and Device}

\textit{Ratified late among those entered in this Log, at the recommendation of the Department of Information Technology, following a matter the Department represented to the Assembly as touching the whole silo's safety. The Assembly's own record of the deliberation preceding the vote does not survive in the Judicial archive in complete form; what is reproduced below is reconstructed from the surviving fragments and the Judge's own summary, entered afterward at Judicial's own insistence that no amendment stand in this Log without some account of its passage.}

``Article 16, Section 6 is amended to add: `Where Information Technology represents to Judicial that a device, signal, or apparatus discovered by a citizen may bear upon matters this Pact does not commit to public record, the citizen who discovers it shall report the discovery to Information Technology directly, and shall not first report it to the Sheriff or to Judicial, notwithstanding any other provision of this Article; and the citizen shall thereafter answer such questions as Information Technology puts to the citizen concerning the manner, place, and circumstance of the discovery. This is not to be understood as a diminishment of Judicial's charge to investigate offenses under this Pact, but as a necessary sequencing of that investigation, that the silo's safety not be made to wait upon a report reaching Information Technology only after having passed through other hands. No citizen who complies with this Section in good faith shall be found to have violated the ordinary reporting requirement this Section otherwise imposes, whatever the outcome of Information Technology's own review proves to be.' The remainder of Article 16 is unchanged.''

\textit{Effect}: A discovery a citizen might once have carried first to the Sheriff or to Judicial is now carried first to Information Technology. The report of the incident that prompted this amendment does not survive in the archive. The concurrence Article 23 requires of the Judge and the Head of Information Technology together, for any amendment touching what this Pact keeps from ordinary record, was given by both offices and is entered in the archive without further note.

\subsection*{Amendment 13: The Narrowing of the Right to Counsel}

\textit{Ratified after Amendment 12, in a season the Assembly's own call to vote described only as necessary to the Pact's continued good order, without further particular given.}

``Article 6, Section 4, as amended by Amendment 6, is further amended: `The right granted by Amendment 6 to bring a speaking-citizen before Judicial does not extend to an offense under Article 16, Sections 1, 3, 5, or 6, it being the judgment of the Assembly that a hearing upon such an offense ought not be made an occasion for the very speech, teaching, relic, or discovery those Sections forbid.' The remainder of Article 6, Section 4 is unchanged.''

\textit{Effect}: A citizen accused of the forbidden word, the false teaching of the founding, the keeping of a restricted instrument, or interference with the apparatus of the silo's keeping may no longer bring a speaking-citizen to the hearing Amendment 6 otherwise grants. The Assembly's record of its own deliberation is entered in the archive in summary only.

\subsection*{Amendment 14: The Recitation Affirmed}

\textit{Ratified after Amendment 13, upon the Mayor's proposal, in a season the Mayor's report called a season of forgetting among the young.}

``Article 15, Section 3 is amended to add: `And the teacher shall, before the recitation is taught, say to the children that the silo is a vessel with a hull, and a hull cannot be pierced in one place without the whole of it taking on water, as the Preamble of this Pact says, that the children know why the recitation is kept, and keep it the more gladly for knowing.' The remainder of Section 3 is unchanged.''

\textit{Effect}: A sentence of the Preamble, which every child already hears read at the opening of the school year under Article 15, Section 9, is to be said again before the recitation is taught. The Judge's statement under Article 21 found no conflict with any provision of this Pact, and no provision altered.

\subsection*{Amendment 15: The Ordinary Pause}

\textit{Ratified some years after Amendment 14, upon the Mayor's proposal at the Sheriff's asking.}

``Article 7, Section 2 is amended to add: `The ordinary pause for rest upon the stair or a landing, beyond which a citizen may not remain there, is a pause of no longer than a count of two hundred, or of four hundred for a citizen carrying a load, or a child, or retired under Article 19; and the Peacekeeper posted at a landing shall keep the count for any citizen who pauses, and move the citizen on at its end.' The remainder of Section 2 is unchanged.''

\textit{Effect}: A pause the Pact left to the Sheriff's judgment is given a length. The Sheriff's report to the Assembly at the sitting following gave the number of citizens moved on under this amendment in its first season as none, the Peacekeepers having found no citizen who paused so long that they could reach the count before the citizen moved of their own accord.

\subsection*{Amendment 16: The Naming of Officers in Assembly}

\textit{Ratified after Amendment 15, upon the Mayor's proposal at the recommendation of the Judge.}

``Article 20, Section 3 is amended to add: `The protection this Section gives a citizen who speaks to the Assembly does not extend to an address that names an officer of the silo and imputes to that officer a breach of this Pact, save where the citizen has first brought the matter to Judicial under Article 6 and Judicial has ruled upon it; an address so made without a ruling is struck from the archive, and the citizen is answerable as for a false bringing under Article 17, Section 5, where Judicial afterward finds the imputation false.' The remainder of Section 3 is unchanged.''

\textit{Effect}: A citizen who would say before the Assembly that an officer has broken this Pact must first have said so to Judicial. The Assembly's record of the sitting enters three addresses struck under this amendment at the sitting next following its ratification, without their text.

\subsection*{Amendment 17: The Mayor's Presence at the Tending}

\textit{Ratified upon a petition of the citizens of Mechanical, the first petition of citizens brought from the Down Deep, after a stopping of the generator for its full tending that ran three days beyond the one day Article 9 allows without the Mayor's written order --- the order having been given, the Head of Mechanical represented to the Assembly, too late to be of use.}

``Article 9, Section 8 is amended to add: `The Mayor shall be present in the Down Deep, at the generator's housing, for the whole of the full tending, and shall sign the log of it beside the Head of Mechanical; and where the tending must run beyond one day, the Mayor's order under this Section is given there, in the log, and not sent.' The remainder of Section 8 is unchanged.''

\textit{Effect}: The Mayor attends the yearly full tending and signs its log. The amendment does not say what the Mayor's presence at the housing is to accomplish, and Mechanical's petition, which survives entire in the archive, does not say either.

\subsection*{Amendment 18: The Notice of Revaluation}

\textit{Ratified twice, the first ratification having been entered in error: the tellers' count under Article 20, Section 6 had added a floor count that arrived after the close of the sitting. The Judge struck the first entry, and the amendment was put again at the sitting following, and carried. Both entries stand in the archive; this Log enters the second.}

``Article 12, Section 3 is amended to add: `No revaluation of chits under this Section takes effect until one cycle after its announcement to the Assembly, and the Chit-Keepers shall post the revaluation, and the day it takes effect, at every market and every dispatch of porters within one day of the announcement.' The remainder of Section 3 is unchanged.''

\textit{Effect}: A revaluation is announced one cycle before it binds. The Assembly's record of the first vote, and the floor count that was wrongly added, are entered under Article 20, Section 6 as that Section provides for a count arriving late.

\subsection*{Amendment 19: The Keeping of the Green List}

\textit{Ratified after Amendment 18, at the recommendation of the Department of Information Technology, upon the Mayor's proposal.}

``Article 22, Section 2 is amended such that the green list is maintained by the Mayor together with the Head of Information Technology, whose concurrence every updating of the list requires, in consultation with Judicial and the Heads of the other Departments as that Section otherwise provides; and the report to the Assembly that the list has been kept is made by the Mayor and the Head of Information Technology together. The remainder of Section 2 is unchanged.''

\textit{Effect}: The green list, which this Pact gave the Mayor to keep in consultation with Judicial and the Heads, is now kept by the Mayor and the Head of Information Technology together, and no updating stands without both. The concurrence Article 23 requires, for an amendment touching what this Pact keeps from ordinary record, was given by the Judge and the Head of Information Technology; the Judge's statement under Article 21 found no conflict with any provision of this Pact.

\subsection*{Amendment 20: The Housing Proviso Restated}

\textit{Ratified last of those entered in this Log, upon a petition of citizens.}

``Article 14, Section 1 is amended to add: `Where a citizen's dwelling is to be revoked under this Section, and the citizen has tried in good faith to mend what occasioned the revoking, and there has been time enough for the mending, the citizen may ask Judicial that the revoking be lessened or stayed, and Judicial may so order.' The remainder of Section 1 is unchanged.''

\textit{Effect}: This Log enters the amendment as ratified. Its effect, beside Amendment 1, is none that the Judge could find, and the Judge's statement under Article 21 says so; the Assembly's record does not show that Amendment 1 was read before the vote.

\section*{Unamended Provisions}

The following provisions of the Pact have not been amended since its founding and remain in their original form:

\begin{itemize}
\item Article 1: Of the Founding and Purpose of the Silo
\item Article 2: Of Citizenship and the Census
\item Article 5: Of the Deputy Mayor and the Order of Succession
\item Article 8: Of the Department of Information Technology
\item Article 10: Of the Department of Mines
\item Article 11: Of Supply, Hydroponics, and the Ration
\item Article 17: Of Crimes and Their Punishments
\item Article 18: Of the Cleaning
\item Article 21: Of Records, Porters, and the Post
\item Article 23: Of Amendment and the Continuance of the Pact
\end{itemize}

The fact that an Article remains unamended does not indicate that it is unchanging or unchangeable---only that no amendment ratified by Assembly has yet altered it. The founding principles of the Pact remain firm; the amendments recorded above show how the silo's citizens have chosen to clarify, extend, restrict, and occasionally soften the Pact's application over time.
